%% kapa: exam=1 round=28 year=2017 period=1 form=A subject=civil_law subject_ko=민법 seq=10 number=10
%% source: Q-Net 감정평가사 기출문제 (cst00309) artlSeq=5199603 file=extracted/1008/5199603/2017년 제28회 감정평가사 1차시험 문제지/1교시 A형.pdf
%% conversion: kordoc 4.24.51 -> tools/gen1.py (auto); review=proofread-sonnet; flags=none
\begin{question}{10}
의사표시에 관한 설명으로 옳지 않은 것은? (다툼이 있으면 판례에 따름)
\begin{choices}
  \item 통정허위표시에서 파산관재인은 제3자에 해당하지 않는다.
  \item 통정허위표시에서 제3자가 보호받기 위해서는 선의이면 되고 그 과실 유무는 묻지 않는다.
  \item 상대방에 의해 유발된 동기의 착오는 동기가 표시되지 않았더라도 법률행위 내용의 중요부분의 착오가 될 수 있다.
  \item 통정허위표시는 제3자 유무와 상관없이 당사자 사이에서는 무효이다.
  \item 사기에 의한 의사표시의 취소는 선의의 제3자에게 대항하지 못한다.
\end{choices}
\end{question}
\answer{1}
